\begin{ouverture}[calcul]{4}{Diviser : partages et restes}{Partie 1 · Nombres entiers et calculs}
{Multiples · Diviseurs · Division euclidienne}
\objectif{Reconnaître multiples et diviseurs avec les critères par 2, 5 et 10.}
\objectif{Poser une division euclidienne avec un diviseur à un chiffre.}
\objectif{Interpréter le quotient et le reste dans un problème.}
\prerequis{Tables de multiplication et soustraction posée.}
\end{ouverture}
\pagestyle{otMain}\essentiel\fontsize{11}{15}\selectfont
\phantomsection\label{cours-4-1}
\essSection{1}{Multiples et diviseurs}
\begin{definition}Comme $36=4\times9$, le nombre 36 est un multiple de 4 et de 9. On dit aussi que 4 et 9 sont des diviseurs de 36 : le partage de 36 en 4 parts égales tombe juste.\end{definition}
\begin{propriete}Un nombre est multiple de 2 si son chiffre des unités est 0, 2, 4, 6 ou 8 : c’est un nombre pair. Il est multiple de 5 si son chiffre des unités est 0 ou 5, et multiple de 10 s’il se termine par 0.\end{propriete}
\begin{exemple}$90$ se termine par 0 : il est multiple de 2, de 5 et de 10. $85$ se termine par 5 : il est multiple de 5, mais pas de 2. Pour savoir si 7 divise 91, aucun critère : on calcule, et $7\times13=91$ convient.\end{exemple}
\begin{savoirfaire}[Tester avec les tables]Pour tester un diviseur de 1 à 10, s’appuyer sur les tables : 54 est dans la table de 6, donc 6 divise 54. Sinon, effectuer la division et regarder si le reste est nul.\end{savoirfaire}
\clearpage
\phantomsection\label{cours-4-2}
\essSection{2}{La division euclidienne}
\begin{definition}Diviser 87 par 4, c’est chercher combien de paquets de 4 remplissent 87, ou la valeur d’une part quand on partage 87 en 4. On écrit $87=4\times21+3$ : le quotient est 21, le reste est 3. Le reste est toujours plus petit que le diviseur.\end{definition}
\begin{savoirfaire}[Poser une division]Pour poser $1\,258:6$ : dans 12, combien de fois 6 ? Deux fois. On abaisse les chiffres un à un. Le quotient est 209 et le reste 4. Quand un chiffre abaissé ne suffit pas, écrire 0 au quotient avant de continuer.\end{savoirfaire}
\begin{exemple}Contrôle d’une division : $6\times209+4=1\,254+4=1\,258$. Le calcul retrouve le nombre de départ, et $4<6$ : la division est juste.\end{exemple}
\begin{attention}La question décide de la réponse. Avec 230 œufs et des boîtes de 6, on remplit 38 boîtes complètes ; mais il faut 39 boîtes pour emporter tous les œufs. Le même reste se lit de deux façons.\end{attention}

\clearpage\rubrique{Découverte}
\begin{blocLivre}{ex-4-1}\begin{exercice}[1]{Les premiers multiples}
Écrire les cinq premiers multiples de 8 non nuls.
\end{exercice}\end{blocLivre}
\begin{blocLivre}{ex-4-2}\begin{exercice}[1]{Pair ou impair}
Trier ces nombres en pairs et impairs : 346 ; $1\,275$ ; 90 ; 83.
\end{exercice}\end{blocLivre}
\begin{blocLivre}{ex-4-3}\begin{exercice}[1]{Multiples de 5 et de 10}
Parmi 85, 120, 905 et $1\,004$, lesquels sont des multiples de 5 ? de 10 ?
\end{exercice}\end{blocLivre}
\begin{blocLivre}{ex-4-4}\begin{exercice}[1]{Une division posée}
Poser la division euclidienne de 87 par 4, puis écrire l’égalité correspondante.
\end{exercice}\end{blocLivre}
\begin{blocLivre}{ex-4-5}\begin{exercice}[1]{Diviseur ou pas ?}
7 est-il un diviseur de 91 ? 6 est-il un diviseur de 75 ?
\end{exercice}\end{blocLivre}

\rubrique{Entraînement}
\begin{blocLivre}{ex-4-6}\begin{exercice}[2]{Partager équitablement}
345 billes sont partagées équitablement entre 5 enfants. Combien de billes reçoit chacun ?
\end{exercice}\end{blocLivre}
\begin{blocLivre}{ex-4-7}\begin{exercice}[2]{Quotient et contrôle}
Effectuer la division euclidienne de $1\,258$ par 6, puis contrôler le résultat.
\end{exercice}\end{blocLivre}
\begin{blocLivre}{ex-4-8}\begin{exercice}[2]{Encadrer par des multiples}
Encadrer 50 entre deux multiples consécutifs de 7.
\end{exercice}\end{blocLivre}
\begin{blocLivre}{ex-4-9}\begin{exercice}[2]{Un reste trop grand}
Lina écrit : $74=8\times8+10$. Expliquer son erreur et corriger l’égalité.
\end{exercice}\end{blocLivre}
\begin{blocLivre}{ex-4-10}\begin{exercice}[2]{Moitiés et quarts}
Calculer la moitié de 470, puis le quart de 232.
\end{exercice}\end{blocLivre}

\rubrique{Problèmes}
\begin{blocLivre}{ex-4-11}\begin{exercice}[3]{Les œufs en boîtes}
Une ferme range 230 œufs dans des boîtes de 6.\begin{questions}
\item Combien de boîtes complètes remplit-elle ?
\item Combien de boîtes faut-il pour emporter tous les œufs ?
\end{questions}
\end{exercice}\end{blocLivre}
\begin{blocLivre}{ex-4-12}\begin{exercice}[3]{Le vide-grenier}
Pour un vide-grenier, 91 livres sont répartis en 7 lots identiques. Combien de livres par lot ? Le partage tombe-t-il juste ?
\end{exercice}\end{blocLivre}
\begin{blocLivre}{ex-4-13}\begin{exercice}[3]{Des équipes de huit}
134 élèves forment des équipes de 8 pour un tournoi.\begin{questions}
\item Combien d’équipes complètes peut-on former ?
\item Combien d’élèves ne sont pas dans une équipe complète ?
\item Combien d’élèves faudrait-il en plus pour une équipe complète de plus ?
\end{questions}
\end{exercice}\end{blocLivre}
\begin{blocLivre}{ex-4-14}\begin{exercice}[3]{Le nombre mystère}
Je suis un multiple de 5, je suis compris entre 82 et 94, et je suis pair. Qui suis-je ? Expliquer pourquoi la réponse est unique.
\end{exercice}\end{blocLivre}
\begin{blocLivre}{ex-4-15}\begin{exercice}[3]{Toujours vrai ?}
Zoé affirme : « La somme de deux nombres pairs est toujours paire. » Tester avec $26+58$ et $134+72$, puis expliquer avec des paquets de deux.
\end{exercice}\end{blocLivre}

\clearpage
